%=========== ΛΥΣΗ - 109ο ΘΕΜΑ Β ===========

\begin{thema}{Β}
Από το σχήμα γνωρίζουμε ότι η $C_f$ έχει οριζόντια ασύμπτωτη την ευθεία $y=2$ στο $-\infty$, μοναδικό ολικό ελάχιστο το $A(0,-1)$ και τέμνει τον άξονα $x'x$ μόνο στα σημεία με τετμημένες $-2$ και $2$.

\begin{erwthma}
\item Επειδή η ευθεία $y=2$ είναι οριζόντια ασύμπτωτη της $C_f$ στο $-\infty$, έχουμε
\[
\lim_{x\to-\infty}f(x)=2.
\]
Άρα, από τη συνέχεια της εκθετικής συνάρτησης,
\[
L_1=\lim_{x\to-\infty}e^{f(x)}=e^2.
\]

Επίσης, επειδή η $f$ παρουσιάζει μοναδικό ολικό ελάχιστο στο $A(0,-1)$, ισχύει
\[
f(x)\geq -1,\quad x\in\mathbb{R},
\]
με ισότητα μόνο για $x=0$. Άρα για $x\neq 0$ έχουμε $f(x)+1>0$ και, καθώς $x\to 0$,
\[
f(x)+1\to f(0)+1=0.
\]
Επομένως
\[
L_2=\lim_{x\to 0}\frac{1}{f(x)+1}=+\infty.
\]

\item Επειδή η $f$ είναι συνεχής και μηδενίζεται μόνο στα $x=-2$ και $x=2$, ενώ $f(0)=-1<0$, από το σχήμα προκύπτει ότι
\[
f(t)<0 \Leftrightarrow -2<t<2.
\]
Θέτουμε
\[
t=f(x)+2.
\]
Τότε
\[
f(f(x)+2)<0
\Leftrightarrow
-2<f(x)+2<2
\Leftrightarrow
-4<f(x)<0.
\]
Όμως το ολικό ελάχιστο της $f$ είναι $-1$, άρα $f(x)\geq -1>-4$ για κάθε $x\in\mathbb{R}$. Επομένως η ανίσωση ισοδυναμεί με
\[
f(x)<0.
\]
Από το σχήμα αυτό συμβαίνει ακριβώς για
\[
-2<x<2.
\]
Άρα η λύση της ανίσωσης είναι
\[
x\in(-2,2).
\]

\item Αφού η $C_f$ δέχεται πλάγια ασύμπτωτη στο $+\infty$ την ευθεία $y=x-2$, έχουμε
\[
\lim_{x\to+\infty}\left(f(x)-(x-2)\right)=0.
\]
Άρα γράφουμε
\[
f(x)=x-2+\varepsilon(x),\quad \varepsilon(x)\to 0.
\]
Τότε
\begin{align*}
\lim_{x\to+\infty}\frac{x f(x)-x^2+5x}{2x+\hm{x}}
&=\lim_{x\to+\infty}
\frac{x(x-2+\varepsilon(x))-x^2+5x}{2x+\hm{x}}\\
&=\lim_{x\to+\infty}
\frac{3x+x\varepsilon(x)}{2x+\hm{x}}\\
&=\lim_{x\to+\infty}
\frac{3+\varepsilon(x)}{2+\frac{\hm{x}}{x}}
=\frac{3}{2}.
\end{align*}

\item Η $f$ είναι συνεχής στο $[0,2]$ και, σύμφωνα με την υπόθεση, παραγωγίσιμη στο $(0,2)$. Επίσης, από το σχήμα έχουμε
\[
f(0)=-1
\quad\text{και}\quad
f(2)=0.
\]
Εφαρμόζουμε το Θεώρημα Μέσης Τιμής στο διάστημα $[0,2]$. Υπάρχει $\xi\in(0,2)$ τέτοιο ώστε
\[
f'(\xi)=\frac{f(2)-f(0)}{2-0}
=\frac{0-(-1)}{2}
=\frac{1}{2}.
\]
Η εφαπτομένη της $C_f$ στο σημείο με τετμημένη $\xi$ έχει συντελεστή διεύθυνσης $f'(\xi)=\frac12>0$. Άρα σχηματίζει με τον άξονα $x'x$ γωνία $\omega\in\left(0,\frac{\pi}{2}\right)$ και
\[
\ef{\omega}=f'(\xi)=\frac{1}{2}.
\]
\end{erwthma}
\end{thema}

%=========== ΤΕΛΟΣ ΛΥΣΗΣ - 109ο ΘΕΜΑ Β ===========
